\documentclass{report}
\usepackage{amsmath,amssymb}
\setlength{\parindent}{0mm}
\setlength{\parskip}{1em}
\begin{document}
\begin{center}
\rule{6in}{1pt} \
{\large Karen E Willcox \\
{\bf Bringing CSE thinking to education }}

77 Massachusetts Avenue \\ Room 33-412a \\ Cambridge \\ MA 02139
\\
{\tt kwillcox@mit.edu}\end{center}

Perhaps the greatest iterative method of all is education. But imagine a
GMRES solver with no structured way to feed back information from one
iteration to the next. Imagine an optimization algorithm that keeps
attempting to iterate with the same failed step. Without structure and
feedback, our methods lack robustness, stability and convergence. Yet
this is largely the case in education. This talk will discuss some of our
recent and ongoing work to develop the concept of educational maps that
relate topics (crosslinks.mit.edu) and learning outcomes (Xoces.mit.edu)
across an engineering undergraduate curriculum. The talk will also
briefly discuss an ongoing project that applies ideas inspired by
feedback control in aerospace engineering to design blended learning
technology.


\end{document}
